\documentclass{amsart}
% For dealing with references we use the comment environment
\usepackage{verbatim}
\newenvironment{reference}{\comment}{\endcomment}
%\newenvironment{reference}{}{}
\newenvironment{slogan}{\comment}{\endcomment}
\newenvironment{history}{\comment}{\endcomment}

% For commutative diagrams we use Xy-pic
\usepackage[all]{xy}

% We use 2cell for 2-commutative diagrams.
\xyoption{2cell}
\UseAllTwocells

% We use multicol for the list of chapters between chapters
\usepackage{multicol}

% This is generally recommended for better output
\usepackage{lmodern}
\usepackage[T1]{fontenc}

% For cross-file-references

% Package for hypertext links:
\usepackage{hyperref}

% For any local file, say "hello.tex" you want to link to please

% Theorem environments.
%
\theoremstyle{plain}
\newtheorem{theorem}[subsection]{Theorem}
\newtheorem{proposition}[subsection]{Proposition}
\newtheorem{lemma}[subsection]{Lemma}

\theoremstyle{definition}
\newtheorem{definition}[subsection]{Definition}
\newtheorem{example}[subsection]{Example}
\newtheorem{exercise}[subsection]{Exercise}
\newtheorem{situation}[subsection]{Situation}

\theoremstyle{remark}
\newtheorem{remark}[subsection]{Remark}
\newtheorem{remarks}[subsection]{Remarks}

\numberwithin{equation}{subsection}

% Macros
%
\def\lim{\mathop{\mathrm{lim}}\nolimits}
\def\colim{\mathop{\mathrm{colim}}\nolimits}
\def\Spec{\mathop{\mathrm{Spec}}}
\def\Hom{\mathop{\mathrm{Hom}}\nolimits}
\def\Ext{\mathop{\mathrm{Ext}}\nolimits}
\def\SheafHom{\mathop{\mathcal{H}\!\mathit{om}}\nolimits}
\def\SheafExt{\mathop{\mathcal{E}\!\mathit{xt}}\nolimits}
\def\Sch{\mathit{Sch}}
\def\Mor{\mathop{\mathrm{Mor}}\nolimits}
\def\Ob{\mathop{\mathrm{Ob}}\nolimits}
\def\Sh{\mathop{\mathit{Sh}}\nolimits}
\def\NL{\mathop{N\!L}\nolimits}
\def\CH{\mathop{\mathrm{CH}}\nolimits}
\def\proetale{{pro\text{-}\acute{e}tale}}
\def\etale{{\acute{e}tale}}
\def\QCoh{\mathit{QCoh}}
\def\Ker{\mathop{\mathrm{Ker}}}
\def\Im{\mathop{\mathrm{Im}}}
\def\Coker{\mathop{\mathrm{Coker}}}
\def\Coim{\mathop{\mathrm{Coim}}}

% Boxtimes
%
\DeclareMathSymbol{\boxtimes}{\mathbin}{AMSa}{"02}

%
% Macros for moduli stacks/spaces
%
\def\QCohstack{\mathcal{QC}\!\mathit{oh}}
\def\Cohstack{\mathcal{C}\!\mathit{oh}}
\def\Spacesstack{\mathcal{S}\!\mathit{paces}}
\def\Quotfunctor{\mathrm{Quot}}
\def\Hilbfunctor{\mathrm{Hilb}}
\def\Curvesstack{\mathcal{C}\!\mathit{urves}}
\def\Polarizedstack{\mathcal{P}\!\mathit{olarized}}
\def\Complexesstack{\mathcal{C}\!\mathit{omplexes}}
% \Pic is the operator that assigns to X its picard group, usage \Pic(X)
% \Picardstack_{X/B} denotes the Picard stack of X over B
% \Picardfunctor_{X/B} denotes the Picard functor of X over B
\def\Pic{\mathop{\mathrm{Pic}}\nolimits}
\def\Picardstack{\mathcal{P}\!\mathit{ic}}
\def\Picardfunctor{\mathrm{Pic}}
\def\Deformationcategory{\mathcal{D}\!\mathit{ef}}

\setlength{\emergencystretch}{3em}
\begin{document}
\title[Stacks Project, Tag 0A27]{586 --- Every integral separated finite-type curve over a field is affine or projective}
\maketitle
\noindent Copyright (C) 2005--2025 Johan de Jong. Stacks Project authors. Source: \href{https://stacks.math.columbia.edu/tag/0A27}{Tag 0A27}.

\noindent Editorial difficulty note (not author text). Difficulty H1: The proof constructs a function with a pole at a single regular boundary point using the growth of infinitesimal divisor sections and minimal pole order. Its two sections give a finite morphism to the projective line whose affine-chart inverse image is the desired punctured curve, reducing all nonproper curves to this construction..

\noindent Editorial provenance note (not author text). History: excerpt from revision \texttt{a0444\allowbreak 6e57e\allowbreak c1fbc\allowbreak 252a8\allowbreak 71afc\allowbreak ec775\allowbreak 2fb28\allowbreak 07b14}, varieties.tex, lines 9411--9493. Reference presentation was adapted; original-author context and core support blocks were retained or added. The mathematical wording is unchanged. Licensed under GNU FDL 1.2 or later; no invariant sections or cover texts. See ../licenses/COPYING. Context blocks reproduce author definitions and supporting results; they are not additional counted problems.

\begin{definition}
\label{divisors-definition-sum-effective-Cartier-divisors}
Let $S$ be a scheme. Given effective Cartier divisors
$D_1$, $D_2$ on $S$ we set $D = D_1 + D_2$ equal to the
closed subscheme of $S$ corresponding to the quasi-coherent
sheaf of ideals
$\mathcal{I}_{D_1}\mathcal{I}_{D_2} \subset \mathcal{O}_S$.
We call this the {\it sum of the effective Cartier divisors
$D_1$ and $D_2$}.
\end{definition}

\begin{definition}
\label{definition-curve}
Let $k$ be a field. A {\it curve} is a variety of dimension $1$ over $k$.
\end{definition}

\begin{definition}
\label{definition-variety}
Let $k$ be a field. A {\it variety} is a scheme $X$ over $k$
such that $X$ is integral and the structure morphism
$X \to \Spec(k)$ is separated and of finite type.
\end{definition}

\begin{lemma}
\label{lemma-curve-affine-projective}
Let $X$ be a curve over $k$. Then either $X$ is an affine scheme or $X$
is H-projective over $k$.
\end{lemma}

\begin{proof}
Choose $X \to \overline{X}$ with
$\overline{X} \setminus X = \{x_1, \ldots, x_r\}$ as in
Lemma \href{https://stacks.math.columbia.edu/tag/0BXW}{lemma reduced dim 1 projective completion (Tag 0BXW)}.
Then $\overline{X}$ is a curve as well.
If $r = 0$, then $X = \overline{X}$ is H-projective over $k$.
Thus we may assume $r \geq 1$ and our goal is to show that
$X$ is affine. By Lemma \href{https://stacks.math.columbia.edu/tag/09N9}{lemma open in affine curve affine (Tag 09N9)}
it suffices to show that $\overline{X} \setminus \{x_1\}$ is affine.
This reduces us to the claim stated in the next paragraph.

\medskip\noindent
Let $X$ be an H-projective curve over $k$. Let $x \in X$ be a closed point
such that $\mathcal{O}_{X, x}$ is a discrete valuation ring. Claim:
$U = X \setminus \{x\}$ is affine. By Lemma \href{https://stacks.math.columbia.edu/tag/0B8Y}{lemma regular point on curve (Tag 0B8Y)}
the point $x$ defines an effective Cartier divisor of $X$.
For $n \geq 1$ denote $nx = x + \ldots + x$ the $n$-fold sum, see
Divisors, Definition \ref{divisors-definition-sum-effective-Cartier-divisors}.
Denote $\mathcal{O}_{nx}$ the structure sheaf of $nx$ viewed as a
coherent module on $X$.
Since every invertible module on the local scheme $nx$ is trivial
the first short exact sequence of
Divisors, Remark \href{https://stacks.math.columbia.edu/tag/0C6K}{remark ses regular section (Tag 0C6K)}
reads
$$
0 \to \mathcal{O}_X
\xrightarrow{1} \mathcal{O}_X(nx) \to \mathcal{O}_{nx} \to 0
$$
in our case. Note that $\dim_k H^0(X, \mathcal{O}_{nx}) \geq n$.
Namely, by Lemma \href{https://stacks.math.columbia.edu/tag/0AYT}{lemma chi tensor finite (Tag 0AYT)} we have
$H^0(X, \mathcal{O}_{nx}) = \mathcal{O}_{X, x}/(\pi^n)$ where $\pi$ in
$\mathcal{O}_{X, x}$ is a uniformizer and the powers $\pi^i$ map to
$k$-linearly independent elements in $\mathcal{O}_{X, x}/(\pi^n)$
for $i = 0, 1, \ldots, n - 1$. We have $\dim_k H^1(X, \mathcal{O}_X) < \infty$
by Cohomology of Schemes, Lemma
\href{https://stacks.math.columbia.edu/tag/02O6}{lemma proper over affine cohomology finite (Tag 02O6)}.
If $n > \dim_k H^1(X, \mathcal{O}_X)$
we conclude from the long exact cohomology sequence
that there exists an $s \in \Gamma(X, \mathcal{O}_X(nx))$
which is not a section of $\mathcal{O}_X$.
If we take $n$ minimal with this property, then $s$ will
map to a generator of the stalk
$\left(\mathcal{O}_X(nx)\right)_x$ since otherwise it would
define a section of $\mathcal{O}_X((n - 1)x) \subset \mathcal{O}_X(nx)$.
For this $n$ we conclude that $s_0 = 1$ and $s_1 = s$
generate the invertible module $\mathcal{L} = \mathcal{O}_X(nx)$.

\medskip\noindent
Consider the corresponding morphism
$f = \varphi_{\mathcal{L}, (s_0, s_1)} : X \to \mathbf{P}^1_k$
of Constructions, Section \href{https://stacks.math.columbia.edu/tag/01ND}{section projective space (Tag 01ND)}.
Observe that the inverse image of $D_{+}(T_0)$ is $U = X \setminus \{x\}$
as the section $s_0$ of $\mathcal{L}$ only vanishes at $x$.
In particular, $f$ is non-constant, i.e., $\Im(f)$ has more
than one point. Hence $f$ must map the generic
point $\eta$ of $X$ to the generic point of $\mathbf{P}^1_k$.
Hence if $y \in \mathbf{P}^1_k$ is a closed point, then $f^{-1}(\{y\})$
is a closed set of $X$ not containing $\eta$, hence finite.
Finally, $f$ is proper\footnote{Namely, a H-projective variety
is a proper variety by Morphisms, Lemma
\href{https://stacks.math.columbia.edu/tag/0BCL}{lemma projective is quasi projective proper (Tag 0BCL)}.
A morphism of varieties whose source is a proper variety is a proper morphism
by Morphisms, Lemma \href{https://stacks.math.columbia.edu/tag/01W6}{lemma image proper scheme closed (Tag 01W6)}.}.
By Cohomology of Schemes, Lemma
\href{https://stacks.math.columbia.edu/tag/02OH}{lemma proper finite fibre finite in neighbourhood (Tag 02OH)}\footnote{One
can avoid using this lemma which relies on the theorem of formal
functions. Namely, $X$ is projective hence it suffices to show
a proper morphism $f : X \to Y$ with finite fibres between quasi-projective
schemes over $k$ is finite. To do this, one chooses an affine open of $X$
containing the fibre of $f$ over a point $y$ using that any finite set of
points of a quasi-projective scheme over $k$ is contained in an affine.
Shrinking $Y$ to a small affine neighbourhood of $y$ one reduces to the
case of a proper morphism between affines. Such a morphism is finite by
Morphisms, Lemma \href{https://stacks.math.columbia.edu/tag/01WM}{lemma integral universally closed (Tag 01WM)}.}
we conclude that $f$ is finite. Hence $U = f^{-1}(D_{+}(T_0))$
is affine.
\end{proof}
\end{document}
